\documentclass[11pt,a4paper]{article}
% ============================================
% GIFT v3.4 — Shared Preamble
% ============================================
% Usage: % ============================================
% GIFT v3.4 — Shared Preamble
% ============================================
% Usage: % ============================================
% GIFT v3.4 — Shared Preamble
% ============================================
% Usage: \input{gift_preamble} after \documentclass

% ENCODING & FONTS
\usepackage[utf8]{inputenc}
\usepackage[T1]{fontenc}
\usepackage{lmodern}

% PAGE LAYOUT (Golden Ratio)
\usepackage[margin=1.618cm, top=2.618cm, bottom=2.618cm]{geometry}

% ESSENTIAL PACKAGES
\usepackage{float}
\usepackage{caption}
\usepackage{subcaption}
\usepackage{setspace}
\usepackage{fancyhdr}
\usepackage{xcolor}
\usepackage{hyperref}
\usepackage{csquotes}
\usepackage{amsmath}
\usepackage{amssymb}
\usepackage{amsthm}
\usepackage{mathtools}
\usepackage{booktabs}
\usepackage{longtable}
\usepackage{array}
\usepackage{multirow}
\usepackage{tikz}
\usepackage{graphicx}
\usepackage{listings}
\usepackage{enumitem}
% Unicode fallbacks (pandoc artifacts that survive markdown→LaTeX)
\DeclareUnicodeCharacter{00B0}{\ensuremath{^\circ}}
\DeclareUnicodeCharacter{221A}{\ensuremath{\sqrt{\,}}}
\DeclareUnicodeCharacter{2713}{\ensuremath{\checkmark}}
\DeclareUnicodeCharacter{2717}{\ensuremath{\times}}
\DeclareUnicodeCharacter{2070}{\ensuremath{^{0}}}
\DeclareUnicodeCharacter{00B9}{\ensuremath{^{1}}}
\DeclareUnicodeCharacter{00B2}{\ensuremath{^{2}}}
\DeclareUnicodeCharacter{00B3}{\ensuremath{^{3}}}
\DeclareUnicodeCharacter{2074}{\ensuremath{^{4}}}
\DeclareUnicodeCharacter{2075}{\ensuremath{^{5}}}
\DeclareUnicodeCharacter{2076}{\ensuremath{^{6}}}
\DeclareUnicodeCharacter{2077}{\ensuremath{^{7}}}
\DeclareUnicodeCharacter{2078}{\ensuremath{^{8}}}
\DeclareUnicodeCharacter{2079}{\ensuremath{^{9}}}
\DeclareUnicodeCharacter{2071}{\ensuremath{^{i}}}
\DeclareUnicodeCharacter{2080}{\ensuremath{_{0}}}
\DeclareUnicodeCharacter{2081}{\ensuremath{_{1}}}
\DeclareUnicodeCharacter{2082}{\ensuremath{_{2}}}
\DeclareUnicodeCharacter{2083}{\ensuremath{_{3}}}
\DeclareUnicodeCharacter{2084}{\ensuremath{_{4}}}
\DeclareUnicodeCharacter{2085}{\ensuremath{_{5}}}
\DeclareUnicodeCharacter{2086}{\ensuremath{_{6}}}
\DeclareUnicodeCharacter{2087}{\ensuremath{_{7}}}
\DeclareUnicodeCharacter{2088}{\ensuremath{_{8}}}
\DeclareUnicodeCharacter{2089}{\ensuremath{_{9}}}
\DeclareUnicodeCharacter{208A}{\ensuremath{_{+}}}
\DeclareUnicodeCharacter{208B}{\ensuremath{_{-}}}
\DeclareUnicodeCharacter{2192}{\ensuremath{\to}}
\DeclareUnicodeCharacter{2502}{|}
\DeclareUnicodeCharacter{25BC}{\ensuremath{\blacktriangledown}}
\DeclareUnicodeCharacter{2500}{\ensuremath{-}}
\DeclareUnicodeCharacter{2514}{\ensuremath{\llcorner}}
\DeclareUnicodeCharacter{251C}{\ensuremath{\vdash}}
\DeclareUnicodeCharacter{2026}{\ldots}
\DeclareUnicodeCharacter{2208}{\ensuremath{\in}}
\DeclareUnicodeCharacter{2229}{\ensuremath{\cap}}
\DeclareUnicodeCharacter{222A}{\ensuremath{\cup}}
\DeclareUnicodeCharacter{2295}{\ensuremath{\oplus}}
\DeclareUnicodeCharacter{2297}{\ensuremath{\otimes}}
\DeclareUnicodeCharacter{2245}{\ensuremath{\cong}}
\DeclareUnicodeCharacter{2248}{\ensuremath{\approx}}
\DeclareUnicodeCharacter{2260}{\ensuremath{\neq}}
\DeclareUnicodeCharacter{2261}{\ensuremath{\equiv}}
\DeclareUnicodeCharacter{2264}{\ensuremath{\leq}}
\DeclareUnicodeCharacter{2265}{\ensuremath{\geq}}
\DeclareUnicodeCharacter{00D7}{\ensuremath{\times}}
\DeclareUnicodeCharacter{00B1}{\ensuremath{\pm}}
\DeclareUnicodeCharacter{2225}{\ensuremath{\|}}
% Greek letters used in body text (pandoc passes them through unchanged)
\DeclareUnicodeCharacter{03B1}{\ensuremath{\alpha}}
\DeclareUnicodeCharacter{03B2}{\ensuremath{\beta}}
\DeclareUnicodeCharacter{03B3}{\ensuremath{\gamma}}
\DeclareUnicodeCharacter{03B4}{\ensuremath{\delta}}
\DeclareUnicodeCharacter{03B5}{\ensuremath{\varepsilon}}
\DeclareUnicodeCharacter{03B6}{\ensuremath{\zeta}}
\DeclareUnicodeCharacter{03B7}{\ensuremath{\eta}}
\DeclareUnicodeCharacter{03B8}{\ensuremath{\theta}}
\DeclareUnicodeCharacter{03B9}{\ensuremath{\iota}}
\DeclareUnicodeCharacter{03BA}{\ensuremath{\kappa}}
\DeclareUnicodeCharacter{03BB}{\ensuremath{\lambda}}
\DeclareUnicodeCharacter{03BC}{\ensuremath{\mu}}
\DeclareUnicodeCharacter{03BD}{\ensuremath{\nu}}
\DeclareUnicodeCharacter{03BE}{\ensuremath{\xi}}
\DeclareUnicodeCharacter{03C0}{\ensuremath{\pi}}
\DeclareUnicodeCharacter{03C1}{\ensuremath{\rho}}
\DeclareUnicodeCharacter{03C3}{\ensuremath{\sigma}}
\DeclareUnicodeCharacter{03C4}{\ensuremath{\tau}}
\DeclareUnicodeCharacter{03C5}{\ensuremath{\upsilon}}
\DeclareUnicodeCharacter{03C6}{\ensuremath{\varphi}}
\DeclareUnicodeCharacter{03C7}{\ensuremath{\chi}}
\DeclareUnicodeCharacter{03C8}{\ensuremath{\psi}}
\DeclareUnicodeCharacter{03C9}{\ensuremath{\omega}}
\DeclareUnicodeCharacter{0394}{\ensuremath{\Delta}}
\DeclareUnicodeCharacter{0398}{\ensuremath{\Theta}}
\DeclareUnicodeCharacter{039B}{\ensuremath{\Lambda}}
\DeclareUnicodeCharacter{03A0}{\ensuremath{\Pi}}
\DeclareUnicodeCharacter{03A3}{\ensuremath{\Sigma}}
\DeclareUnicodeCharacter{03A6}{\ensuremath{\Phi}}
\DeclareUnicodeCharacter{03A8}{\ensuremath{\Psi}}
\DeclareUnicodeCharacter{03A9}{\ensuremath{\Omega}}
\DeclareUnicodeCharacter{2200}{\ensuremath{\forall}}
\DeclareUnicodeCharacter{2203}{\ensuremath{\exists}}
\DeclareUnicodeCharacter{2207}{\ensuremath{\nabla}}
\DeclareUnicodeCharacter{2202}{\ensuremath{\partial}}
\DeclareUnicodeCharacter{2218}{\ensuremath{\circ}}
\DeclareUnicodeCharacter{2032}{\ensuremath{\prime}}
\DeclareUnicodeCharacter{2113}{\ensuremath{\ell}}
\DeclareUnicodeCharacter{210F}{\ensuremath{\hbar}}
\DeclareUnicodeCharacter{27E8}{\ensuremath{\langle}}
\DeclareUnicodeCharacter{27E9}{\ensuremath{\rangle}}
\DeclareUnicodeCharacter{2308}{\ensuremath{\lceil}}
\DeclareUnicodeCharacter{2309}{\ensuremath{\rceil}}
\DeclareUnicodeCharacter{230A}{\ensuremath{\lfloor}}
\DeclareUnicodeCharacter{230B}{\ensuremath{\rfloor}}
\DeclareUnicodeCharacter{222B}{\ensuremath{\int}}
\DeclareUnicodeCharacter{2211}{\ensuremath{\sum}}
\DeclareUnicodeCharacter{220F}{\ensuremath{\prod}}
\DeclareUnicodeCharacter{221E}{\ensuremath{\infty}}
\DeclareUnicodeCharacter{00B5}{\ensuremath{\mu}}
\DeclareUnicodeCharacter{2236}{\ensuremath{\colon}}
\DeclareUnicodeCharacter{2192}{\ensuremath{\to}}
\DeclareUnicodeCharacter{21A6}{\ensuremath{\mapsto}}
\DeclareUnicodeCharacter{21D2}{\ensuremath{\Rightarrow}}
\DeclareUnicodeCharacter{21D4}{\ensuremath{\Leftrightarrow}}
\DeclareUnicodeCharacter{2282}{\ensuremath{\subset}}
\DeclareUnicodeCharacter{2286}{\ensuremath{\subseteq}}
\DeclareUnicodeCharacter{2287}{\ensuremath{\supseteq}}
\DeclareUnicodeCharacter{222B}{\ensuremath{\int}}
\DeclareUnicodeCharacter{2205}{\ensuremath{\emptyset}}
\DeclareUnicodeCharacter{2299}{\ensuremath{\odot}}
\DeclareUnicodeCharacter{2220}{\ensuremath{\angle}}
\DeclareUnicodeCharacter{2135}{\ensuremath{\aleph}}
\DeclareUnicodeCharacter{207B}{\ensuremath{^{-}}}
\DeclareUnicodeCharacter{207A}{\ensuremath{^{+}}}
\DeclareUnicodeCharacter{220E}{\ensuremath{\blacksquare}}
\DeclareUnicodeCharacter{2061}{}
\DeclareUnicodeCharacter{2062}{}
\DeclareUnicodeCharacter{2063}{}
\DeclareUnicodeCharacter{2009}{\,}
\DeclareUnicodeCharacter{200A}{\,}
\DeclareUnicodeCharacter{200B}{}
\DeclareUnicodeCharacter{2212}{\ensuremath{-}}
\DeclareUnicodeCharacter{2194}{\ensuremath{\leftrightarrow}}
\DeclareUnicodeCharacter{2099}{\ensuremath{_{n}}}
\DeclareUnicodeCharacter{2098}{\ensuremath{_{m}}}
\DeclareUnicodeCharacter{2096}{\ensuremath{_{k}}}
\DeclareUnicodeCharacter{2095}{\ensuremath{_{h}}}
\DeclareUnicodeCharacter{2090}{\ensuremath{_{a}}}
\DeclareUnicodeCharacter{2091}{\ensuremath{_{e}}}
\DeclareUnicodeCharacter{2093}{\ensuremath{_{x}}}
\DeclareUnicodeCharacter{1D62}{\ensuremath{_{i}}}
\DeclareUnicodeCharacter{2C7C}{\ensuremath{_{j}}}
\DeclareUnicodeCharacter{2C7B}{\ensuremath{^{e}}}
\DeclareUnicodeCharacter{1D52}{\ensuremath{^{o}}}
\DeclareUnicodeCharacter{1D5}{\ensuremath{^{u}}}
\DeclareUnicodeCharacter{02E1}{\ensuremath{^{l}}}
\DeclareUnicodeCharacter{2295}{\ensuremath{\oplus}}
\DeclareUnicodeCharacter{27FA}{\ensuremath{\Longleftrightarrow}}
\DeclareUnicodeCharacter{27F9}{\ensuremath{\Longrightarrow}}
\DeclareUnicodeCharacter{2299}{\ensuremath{\odot}}
% Provide \tightlist used by pandoc-generated lists (no-op)
\providecommand{\tightlist}{\setlength{\itemsep}{0pt}\setlength{\parskip}{0pt}}
% pandoc longtable column specs: provide \real as passthrough and a
% `none` counter so that calc's \real{x} expansion path works.
\newcounter{none}
\providecommand{\real}[1]{#1}

% LISTINGS
\lstset{
    basicstyle=\small\ttfamily,
    breaklines=true, frame=single,
    keepspaces=true, showstringspaces=false,
    aboveskip=0.8em, belowskip=0.8em
}

% HEADER/FOOTER
\setlength{\headheight}{14pt}
\pagestyle{fancy}
\fancyhf{}
\fancyhead[R]{\thepage}
\renewcommand{\headrulewidth}{0.2pt}

% HYPERREF
\hypersetup{
    colorlinks=true,
    linkcolor=blue, citecolor=blue, urlcolor=blue,
    pdfauthor={Brieuc de La Fourniere}
}

% SPACING
\setstretch{1.2}
\setlength{\parskip}{0.4em}
\setlength{\parindent}{0pt}

% TITLE FORMATTING
\usepackage{titling}
\pretitle{\LARGE\bfseries}
\posttitle{\vspace{-0.4em}}
\preauthor{}
\postauthor{}
\predate{}
\postdate{}
\setlength{\droptitle}{-2.0em}

% CUSTOM COMMANDS — GIFT notation
\newcommand{\E}{\mathrm{E}}
\newcommand{\Gtwo}{\mathrm{G}_2}
\newcommand{\Kseven}{K_7}
\newcommand{\AdS}{\mathrm{AdS}}
\newcommand{\SU}{\mathrm{SU}}
\newcommand{\SO}{\mathrm{SO}}
\newcommand{\U}{\mathrm{U}}
\newcommand{\Sp}{\mathrm{Sp}}
\newcommand{\PSL}{\mathrm{PSL}}
\newcommand{\SM}{\mathrm{SM}}
\newcommand{\Tr}{\mathrm{Tr}}
\newcommand{\rank}{\mathrm{rank}}
\newcommand{\Vol}{\mathrm{Vol}}
\newcommand{\Hol}{\mathrm{Hol}}
\newcommand{\Ric}{\mathrm{Ric}}
\newcommand{\Rm}{\mathrm{Rm}}
\newcommand{\Scal}{\mathrm{Scal}}
\DeclareMathOperator{\diag}{diag}
\DeclareMathOperator{\cond}{cond}
\newcommand{\GIFT}{\textrm{GIFT}}
\newcommand{\CP}{\mathrm{CP}}
\newcommand{\CKM}{\mathrm{CKM}}
\newcommand{\PMNS}{\mathrm{PMNS}}
\newcommand{\EW}{\mathrm{EW}}
\newcommand{\Pl}{\mathrm{Pl}}
\newcommand{\DE}{\mathrm{DE}}
\newcommand{\DM}{\mathrm{DM}}
\newcommand{\KK}{\mathrm{KK}}
\newcommand{\NK}{\mathrm{NK}}
\newcommand{\TCS}{\mathrm{TCS}}
\newcommand{\PINN}{\mathrm{PINN}}
\newcommand{\btwo}{b_2}
\newcommand{\bthree}{b_3}
\newcommand{\typeI}{\textsc{I}}
\newcommand{\typeII}{\textsc{II}}
\newcommand{\typeIII}{\textsc{III}}
\newcommand{\typeIV}{\textsc{IV}}
\newcommand{\proven}{\textsc{Proven}}

% PDF bookmark safety
\pdfstringdefDisableCommands{%
  \def\textsubscript#1{#1}%
  \def\textsuperscript#1{#1}%
  \def\CP{CP}%
  \def\Gtwo{G2}%
  \def\Kseven{K7}%
  \def\GIFT{GIFT}%
  \def\E{E}%
  \def\SM{SM}%
}

% THEOREM ENVIRONMENTS
\newtheorem{theorem}{Theorem}[section]
\newtheorem{proposition}[theorem]{Proposition}
\newtheorem{lemma}[theorem]{Lemma}
\newtheorem{corollary}[theorem]{Corollary}
\theoremstyle{definition}
\newtheorem{definition}[theorem]{Definition}
\theoremstyle{remark}
\newtheorem{remark}[theorem]{Remark}
 after \documentclass

% ENCODING & FONTS
\usepackage[utf8]{inputenc}
\usepackage[T1]{fontenc}
\usepackage{lmodern}

% PAGE LAYOUT (Golden Ratio)
\usepackage[margin=1.618cm, top=2.618cm, bottom=2.618cm]{geometry}

% ESSENTIAL PACKAGES
\usepackage{float}
\usepackage{caption}
\usepackage{subcaption}
\usepackage{setspace}
\usepackage{fancyhdr}
\usepackage{xcolor}
\usepackage{hyperref}
\usepackage{csquotes}
\usepackage{amsmath}
\usepackage{amssymb}
\usepackage{amsthm}
\usepackage{mathtools}
\usepackage{booktabs}
\usepackage{longtable}
\usepackage{array}
\usepackage{multirow}
\usepackage{tikz}
\usepackage{graphicx}
\usepackage{listings}
\usepackage{enumitem}
% Unicode fallbacks (pandoc artifacts that survive markdown→LaTeX)
\DeclareUnicodeCharacter{00B0}{\ensuremath{^\circ}}
\DeclareUnicodeCharacter{221A}{\ensuremath{\sqrt{\,}}}
\DeclareUnicodeCharacter{2713}{\ensuremath{\checkmark}}
\DeclareUnicodeCharacter{2717}{\ensuremath{\times}}
\DeclareUnicodeCharacter{2070}{\ensuremath{^{0}}}
\DeclareUnicodeCharacter{00B9}{\ensuremath{^{1}}}
\DeclareUnicodeCharacter{00B2}{\ensuremath{^{2}}}
\DeclareUnicodeCharacter{00B3}{\ensuremath{^{3}}}
\DeclareUnicodeCharacter{2074}{\ensuremath{^{4}}}
\DeclareUnicodeCharacter{2075}{\ensuremath{^{5}}}
\DeclareUnicodeCharacter{2076}{\ensuremath{^{6}}}
\DeclareUnicodeCharacter{2077}{\ensuremath{^{7}}}
\DeclareUnicodeCharacter{2078}{\ensuremath{^{8}}}
\DeclareUnicodeCharacter{2079}{\ensuremath{^{9}}}
\DeclareUnicodeCharacter{2071}{\ensuremath{^{i}}}
\DeclareUnicodeCharacter{2080}{\ensuremath{_{0}}}
\DeclareUnicodeCharacter{2081}{\ensuremath{_{1}}}
\DeclareUnicodeCharacter{2082}{\ensuremath{_{2}}}
\DeclareUnicodeCharacter{2083}{\ensuremath{_{3}}}
\DeclareUnicodeCharacter{2084}{\ensuremath{_{4}}}
\DeclareUnicodeCharacter{2085}{\ensuremath{_{5}}}
\DeclareUnicodeCharacter{2086}{\ensuremath{_{6}}}
\DeclareUnicodeCharacter{2087}{\ensuremath{_{7}}}
\DeclareUnicodeCharacter{2088}{\ensuremath{_{8}}}
\DeclareUnicodeCharacter{2089}{\ensuremath{_{9}}}
\DeclareUnicodeCharacter{208A}{\ensuremath{_{+}}}
\DeclareUnicodeCharacter{208B}{\ensuremath{_{-}}}
\DeclareUnicodeCharacter{2192}{\ensuremath{\to}}
\DeclareUnicodeCharacter{2502}{|}
\DeclareUnicodeCharacter{25BC}{\ensuremath{\blacktriangledown}}
\DeclareUnicodeCharacter{2500}{\ensuremath{-}}
\DeclareUnicodeCharacter{2514}{\ensuremath{\llcorner}}
\DeclareUnicodeCharacter{251C}{\ensuremath{\vdash}}
\DeclareUnicodeCharacter{2026}{\ldots}
\DeclareUnicodeCharacter{2208}{\ensuremath{\in}}
\DeclareUnicodeCharacter{2229}{\ensuremath{\cap}}
\DeclareUnicodeCharacter{222A}{\ensuremath{\cup}}
\DeclareUnicodeCharacter{2295}{\ensuremath{\oplus}}
\DeclareUnicodeCharacter{2297}{\ensuremath{\otimes}}
\DeclareUnicodeCharacter{2245}{\ensuremath{\cong}}
\DeclareUnicodeCharacter{2248}{\ensuremath{\approx}}
\DeclareUnicodeCharacter{2260}{\ensuremath{\neq}}
\DeclareUnicodeCharacter{2261}{\ensuremath{\equiv}}
\DeclareUnicodeCharacter{2264}{\ensuremath{\leq}}
\DeclareUnicodeCharacter{2265}{\ensuremath{\geq}}
\DeclareUnicodeCharacter{00D7}{\ensuremath{\times}}
\DeclareUnicodeCharacter{00B1}{\ensuremath{\pm}}
\DeclareUnicodeCharacter{2225}{\ensuremath{\|}}
% Greek letters used in body text (pandoc passes them through unchanged)
\DeclareUnicodeCharacter{03B1}{\ensuremath{\alpha}}
\DeclareUnicodeCharacter{03B2}{\ensuremath{\beta}}
\DeclareUnicodeCharacter{03B3}{\ensuremath{\gamma}}
\DeclareUnicodeCharacter{03B4}{\ensuremath{\delta}}
\DeclareUnicodeCharacter{03B5}{\ensuremath{\varepsilon}}
\DeclareUnicodeCharacter{03B6}{\ensuremath{\zeta}}
\DeclareUnicodeCharacter{03B7}{\ensuremath{\eta}}
\DeclareUnicodeCharacter{03B8}{\ensuremath{\theta}}
\DeclareUnicodeCharacter{03B9}{\ensuremath{\iota}}
\DeclareUnicodeCharacter{03BA}{\ensuremath{\kappa}}
\DeclareUnicodeCharacter{03BB}{\ensuremath{\lambda}}
\DeclareUnicodeCharacter{03BC}{\ensuremath{\mu}}
\DeclareUnicodeCharacter{03BD}{\ensuremath{\nu}}
\DeclareUnicodeCharacter{03BE}{\ensuremath{\xi}}
\DeclareUnicodeCharacter{03C0}{\ensuremath{\pi}}
\DeclareUnicodeCharacter{03C1}{\ensuremath{\rho}}
\DeclareUnicodeCharacter{03C3}{\ensuremath{\sigma}}
\DeclareUnicodeCharacter{03C4}{\ensuremath{\tau}}
\DeclareUnicodeCharacter{03C5}{\ensuremath{\upsilon}}
\DeclareUnicodeCharacter{03C6}{\ensuremath{\varphi}}
\DeclareUnicodeCharacter{03C7}{\ensuremath{\chi}}
\DeclareUnicodeCharacter{03C8}{\ensuremath{\psi}}
\DeclareUnicodeCharacter{03C9}{\ensuremath{\omega}}
\DeclareUnicodeCharacter{0394}{\ensuremath{\Delta}}
\DeclareUnicodeCharacter{0398}{\ensuremath{\Theta}}
\DeclareUnicodeCharacter{039B}{\ensuremath{\Lambda}}
\DeclareUnicodeCharacter{03A0}{\ensuremath{\Pi}}
\DeclareUnicodeCharacter{03A3}{\ensuremath{\Sigma}}
\DeclareUnicodeCharacter{03A6}{\ensuremath{\Phi}}
\DeclareUnicodeCharacter{03A8}{\ensuremath{\Psi}}
\DeclareUnicodeCharacter{03A9}{\ensuremath{\Omega}}
\DeclareUnicodeCharacter{2200}{\ensuremath{\forall}}
\DeclareUnicodeCharacter{2203}{\ensuremath{\exists}}
\DeclareUnicodeCharacter{2207}{\ensuremath{\nabla}}
\DeclareUnicodeCharacter{2202}{\ensuremath{\partial}}
\DeclareUnicodeCharacter{2218}{\ensuremath{\circ}}
\DeclareUnicodeCharacter{2032}{\ensuremath{\prime}}
\DeclareUnicodeCharacter{2113}{\ensuremath{\ell}}
\DeclareUnicodeCharacter{210F}{\ensuremath{\hbar}}
\DeclareUnicodeCharacter{27E8}{\ensuremath{\langle}}
\DeclareUnicodeCharacter{27E9}{\ensuremath{\rangle}}
\DeclareUnicodeCharacter{2308}{\ensuremath{\lceil}}
\DeclareUnicodeCharacter{2309}{\ensuremath{\rceil}}
\DeclareUnicodeCharacter{230A}{\ensuremath{\lfloor}}
\DeclareUnicodeCharacter{230B}{\ensuremath{\rfloor}}
\DeclareUnicodeCharacter{222B}{\ensuremath{\int}}
\DeclareUnicodeCharacter{2211}{\ensuremath{\sum}}
\DeclareUnicodeCharacter{220F}{\ensuremath{\prod}}
\DeclareUnicodeCharacter{221E}{\ensuremath{\infty}}
\DeclareUnicodeCharacter{00B5}{\ensuremath{\mu}}
\DeclareUnicodeCharacter{2236}{\ensuremath{\colon}}
\DeclareUnicodeCharacter{2192}{\ensuremath{\to}}
\DeclareUnicodeCharacter{21A6}{\ensuremath{\mapsto}}
\DeclareUnicodeCharacter{21D2}{\ensuremath{\Rightarrow}}
\DeclareUnicodeCharacter{21D4}{\ensuremath{\Leftrightarrow}}
\DeclareUnicodeCharacter{2282}{\ensuremath{\subset}}
\DeclareUnicodeCharacter{2286}{\ensuremath{\subseteq}}
\DeclareUnicodeCharacter{2287}{\ensuremath{\supseteq}}
\DeclareUnicodeCharacter{222B}{\ensuremath{\int}}
\DeclareUnicodeCharacter{2205}{\ensuremath{\emptyset}}
\DeclareUnicodeCharacter{2299}{\ensuremath{\odot}}
\DeclareUnicodeCharacter{2220}{\ensuremath{\angle}}
\DeclareUnicodeCharacter{2135}{\ensuremath{\aleph}}
\DeclareUnicodeCharacter{207B}{\ensuremath{^{-}}}
\DeclareUnicodeCharacter{207A}{\ensuremath{^{+}}}
\DeclareUnicodeCharacter{220E}{\ensuremath{\blacksquare}}
\DeclareUnicodeCharacter{2061}{}
\DeclareUnicodeCharacter{2062}{}
\DeclareUnicodeCharacter{2063}{}
\DeclareUnicodeCharacter{2009}{\,}
\DeclareUnicodeCharacter{200A}{\,}
\DeclareUnicodeCharacter{200B}{}
\DeclareUnicodeCharacter{2212}{\ensuremath{-}}
\DeclareUnicodeCharacter{2194}{\ensuremath{\leftrightarrow}}
\DeclareUnicodeCharacter{2099}{\ensuremath{_{n}}}
\DeclareUnicodeCharacter{2098}{\ensuremath{_{m}}}
\DeclareUnicodeCharacter{2096}{\ensuremath{_{k}}}
\DeclareUnicodeCharacter{2095}{\ensuremath{_{h}}}
\DeclareUnicodeCharacter{2090}{\ensuremath{_{a}}}
\DeclareUnicodeCharacter{2091}{\ensuremath{_{e}}}
\DeclareUnicodeCharacter{2093}{\ensuremath{_{x}}}
\DeclareUnicodeCharacter{1D62}{\ensuremath{_{i}}}
\DeclareUnicodeCharacter{2C7C}{\ensuremath{_{j}}}
\DeclareUnicodeCharacter{2C7B}{\ensuremath{^{e}}}
\DeclareUnicodeCharacter{1D52}{\ensuremath{^{o}}}
\DeclareUnicodeCharacter{1D5}{\ensuremath{^{u}}}
\DeclareUnicodeCharacter{02E1}{\ensuremath{^{l}}}
\DeclareUnicodeCharacter{2295}{\ensuremath{\oplus}}
\DeclareUnicodeCharacter{27FA}{\ensuremath{\Longleftrightarrow}}
\DeclareUnicodeCharacter{27F9}{\ensuremath{\Longrightarrow}}
\DeclareUnicodeCharacter{2299}{\ensuremath{\odot}}
% Provide \tightlist used by pandoc-generated lists (no-op)
\providecommand{\tightlist}{\setlength{\itemsep}{0pt}\setlength{\parskip}{0pt}}
% pandoc longtable column specs: provide \real as passthrough and a
% `none` counter so that calc's \real{x} expansion path works.
\newcounter{none}
\providecommand{\real}[1]{#1}

% LISTINGS
\lstset{
    basicstyle=\small\ttfamily,
    breaklines=true, frame=single,
    keepspaces=true, showstringspaces=false,
    aboveskip=0.8em, belowskip=0.8em
}

% HEADER/FOOTER
\setlength{\headheight}{14pt}
\pagestyle{fancy}
\fancyhf{}
\fancyhead[R]{\thepage}
\renewcommand{\headrulewidth}{0.2pt}

% HYPERREF
\hypersetup{
    colorlinks=true,
    linkcolor=blue, citecolor=blue, urlcolor=blue,
    pdfauthor={Brieuc de La Fourniere}
}

% SPACING
\setstretch{1.2}
\setlength{\parskip}{0.4em}
\setlength{\parindent}{0pt}

% TITLE FORMATTING
\usepackage{titling}
\pretitle{\LARGE\bfseries}
\posttitle{\vspace{-0.4em}}
\preauthor{}
\postauthor{}
\predate{}
\postdate{}
\setlength{\droptitle}{-2.0em}

% CUSTOM COMMANDS — GIFT notation
\newcommand{\E}{\mathrm{E}}
\newcommand{\Gtwo}{\mathrm{G}_2}
\newcommand{\Kseven}{K_7}
\newcommand{\AdS}{\mathrm{AdS}}
\newcommand{\SU}{\mathrm{SU}}
\newcommand{\SO}{\mathrm{SO}}
\newcommand{\U}{\mathrm{U}}
\newcommand{\Sp}{\mathrm{Sp}}
\newcommand{\PSL}{\mathrm{PSL}}
\newcommand{\SM}{\mathrm{SM}}
\newcommand{\Tr}{\mathrm{Tr}}
\newcommand{\rank}{\mathrm{rank}}
\newcommand{\Vol}{\mathrm{Vol}}
\newcommand{\Hol}{\mathrm{Hol}}
\newcommand{\Ric}{\mathrm{Ric}}
\newcommand{\Rm}{\mathrm{Rm}}
\newcommand{\Scal}{\mathrm{Scal}}
\DeclareMathOperator{\diag}{diag}
\DeclareMathOperator{\cond}{cond}
\newcommand{\GIFT}{\textrm{GIFT}}
\newcommand{\CP}{\mathrm{CP}}
\newcommand{\CKM}{\mathrm{CKM}}
\newcommand{\PMNS}{\mathrm{PMNS}}
\newcommand{\EW}{\mathrm{EW}}
\newcommand{\Pl}{\mathrm{Pl}}
\newcommand{\DE}{\mathrm{DE}}
\newcommand{\DM}{\mathrm{DM}}
\newcommand{\KK}{\mathrm{KK}}
\newcommand{\NK}{\mathrm{NK}}
\newcommand{\TCS}{\mathrm{TCS}}
\newcommand{\PINN}{\mathrm{PINN}}
\newcommand{\btwo}{b_2}
\newcommand{\bthree}{b_3}
\newcommand{\typeI}{\textsc{I}}
\newcommand{\typeII}{\textsc{II}}
\newcommand{\typeIII}{\textsc{III}}
\newcommand{\typeIV}{\textsc{IV}}
\newcommand{\proven}{\textsc{Proven}}

% PDF bookmark safety
\pdfstringdefDisableCommands{%
  \def\textsubscript#1{#1}%
  \def\textsuperscript#1{#1}%
  \def\CP{CP}%
  \def\Gtwo{G2}%
  \def\Kseven{K7}%
  \def\GIFT{GIFT}%
  \def\E{E}%
  \def\SM{SM}%
}

% THEOREM ENVIRONMENTS
\newtheorem{theorem}{Theorem}[section]
\newtheorem{proposition}[theorem]{Proposition}
\newtheorem{lemma}[theorem]{Lemma}
\newtheorem{corollary}[theorem]{Corollary}
\theoremstyle{definition}
\newtheorem{definition}[theorem]{Definition}
\theoremstyle{remark}
\newtheorem{remark}[theorem]{Remark}
 after \documentclass

% ENCODING & FONTS
\usepackage[utf8]{inputenc}
\usepackage[T1]{fontenc}
\usepackage{lmodern}

% PAGE LAYOUT (Golden Ratio)
\usepackage[margin=1.618cm, top=2.618cm, bottom=2.618cm]{geometry}

% ESSENTIAL PACKAGES
\usepackage{float}
\usepackage{caption}
\usepackage{subcaption}
\usepackage{setspace}
\usepackage{fancyhdr}
\usepackage{xcolor}
\usepackage{hyperref}
\usepackage{csquotes}
\usepackage{amsmath}
\usepackage{amssymb}
\usepackage{amsthm}
\usepackage{mathtools}
\usepackage{booktabs}
\usepackage{longtable}
\usepackage{array}
\usepackage{multirow}
\usepackage{tikz}
\usepackage{graphicx}
\usepackage{listings}
\usepackage{enumitem}
% Unicode fallbacks (pandoc artifacts that survive markdown→LaTeX)
\DeclareUnicodeCharacter{00B0}{\ensuremath{^\circ}}
\DeclareUnicodeCharacter{221A}{\ensuremath{\sqrt{\,}}}
\DeclareUnicodeCharacter{2713}{\ensuremath{\checkmark}}
\DeclareUnicodeCharacter{2717}{\ensuremath{\times}}
\DeclareUnicodeCharacter{2070}{\ensuremath{^{0}}}
\DeclareUnicodeCharacter{00B9}{\ensuremath{^{1}}}
\DeclareUnicodeCharacter{00B2}{\ensuremath{^{2}}}
\DeclareUnicodeCharacter{00B3}{\ensuremath{^{3}}}
\DeclareUnicodeCharacter{2074}{\ensuremath{^{4}}}
\DeclareUnicodeCharacter{2075}{\ensuremath{^{5}}}
\DeclareUnicodeCharacter{2076}{\ensuremath{^{6}}}
\DeclareUnicodeCharacter{2077}{\ensuremath{^{7}}}
\DeclareUnicodeCharacter{2078}{\ensuremath{^{8}}}
\DeclareUnicodeCharacter{2079}{\ensuremath{^{9}}}
\DeclareUnicodeCharacter{2071}{\ensuremath{^{i}}}
\DeclareUnicodeCharacter{2080}{\ensuremath{_{0}}}
\DeclareUnicodeCharacter{2081}{\ensuremath{_{1}}}
\DeclareUnicodeCharacter{2082}{\ensuremath{_{2}}}
\DeclareUnicodeCharacter{2083}{\ensuremath{_{3}}}
\DeclareUnicodeCharacter{2084}{\ensuremath{_{4}}}
\DeclareUnicodeCharacter{2085}{\ensuremath{_{5}}}
\DeclareUnicodeCharacter{2086}{\ensuremath{_{6}}}
\DeclareUnicodeCharacter{2087}{\ensuremath{_{7}}}
\DeclareUnicodeCharacter{2088}{\ensuremath{_{8}}}
\DeclareUnicodeCharacter{2089}{\ensuremath{_{9}}}
\DeclareUnicodeCharacter{208A}{\ensuremath{_{+}}}
\DeclareUnicodeCharacter{208B}{\ensuremath{_{-}}}
\DeclareUnicodeCharacter{2192}{\ensuremath{\to}}
\DeclareUnicodeCharacter{2502}{|}
\DeclareUnicodeCharacter{25BC}{\ensuremath{\blacktriangledown}}
\DeclareUnicodeCharacter{2500}{\ensuremath{-}}
\DeclareUnicodeCharacter{2514}{\ensuremath{\llcorner}}
\DeclareUnicodeCharacter{251C}{\ensuremath{\vdash}}
\DeclareUnicodeCharacter{2026}{\ldots}
\DeclareUnicodeCharacter{2208}{\ensuremath{\in}}
\DeclareUnicodeCharacter{2229}{\ensuremath{\cap}}
\DeclareUnicodeCharacter{222A}{\ensuremath{\cup}}
\DeclareUnicodeCharacter{2295}{\ensuremath{\oplus}}
\DeclareUnicodeCharacter{2297}{\ensuremath{\otimes}}
\DeclareUnicodeCharacter{2245}{\ensuremath{\cong}}
\DeclareUnicodeCharacter{2248}{\ensuremath{\approx}}
\DeclareUnicodeCharacter{2260}{\ensuremath{\neq}}
\DeclareUnicodeCharacter{2261}{\ensuremath{\equiv}}
\DeclareUnicodeCharacter{2264}{\ensuremath{\leq}}
\DeclareUnicodeCharacter{2265}{\ensuremath{\geq}}
\DeclareUnicodeCharacter{00D7}{\ensuremath{\times}}
\DeclareUnicodeCharacter{00B1}{\ensuremath{\pm}}
\DeclareUnicodeCharacter{2225}{\ensuremath{\|}}
% Greek letters used in body text (pandoc passes them through unchanged)
\DeclareUnicodeCharacter{03B1}{\ensuremath{\alpha}}
\DeclareUnicodeCharacter{03B2}{\ensuremath{\beta}}
\DeclareUnicodeCharacter{03B3}{\ensuremath{\gamma}}
\DeclareUnicodeCharacter{03B4}{\ensuremath{\delta}}
\DeclareUnicodeCharacter{03B5}{\ensuremath{\varepsilon}}
\DeclareUnicodeCharacter{03B6}{\ensuremath{\zeta}}
\DeclareUnicodeCharacter{03B7}{\ensuremath{\eta}}
\DeclareUnicodeCharacter{03B8}{\ensuremath{\theta}}
\DeclareUnicodeCharacter{03B9}{\ensuremath{\iota}}
\DeclareUnicodeCharacter{03BA}{\ensuremath{\kappa}}
\DeclareUnicodeCharacter{03BB}{\ensuremath{\lambda}}
\DeclareUnicodeCharacter{03BC}{\ensuremath{\mu}}
\DeclareUnicodeCharacter{03BD}{\ensuremath{\nu}}
\DeclareUnicodeCharacter{03BE}{\ensuremath{\xi}}
\DeclareUnicodeCharacter{03C0}{\ensuremath{\pi}}
\DeclareUnicodeCharacter{03C1}{\ensuremath{\rho}}
\DeclareUnicodeCharacter{03C3}{\ensuremath{\sigma}}
\DeclareUnicodeCharacter{03C4}{\ensuremath{\tau}}
\DeclareUnicodeCharacter{03C5}{\ensuremath{\upsilon}}
\DeclareUnicodeCharacter{03C6}{\ensuremath{\varphi}}
\DeclareUnicodeCharacter{03C7}{\ensuremath{\chi}}
\DeclareUnicodeCharacter{03C8}{\ensuremath{\psi}}
\DeclareUnicodeCharacter{03C9}{\ensuremath{\omega}}
\DeclareUnicodeCharacter{0394}{\ensuremath{\Delta}}
\DeclareUnicodeCharacter{0398}{\ensuremath{\Theta}}
\DeclareUnicodeCharacter{039B}{\ensuremath{\Lambda}}
\DeclareUnicodeCharacter{03A0}{\ensuremath{\Pi}}
\DeclareUnicodeCharacter{03A3}{\ensuremath{\Sigma}}
\DeclareUnicodeCharacter{03A6}{\ensuremath{\Phi}}
\DeclareUnicodeCharacter{03A8}{\ensuremath{\Psi}}
\DeclareUnicodeCharacter{03A9}{\ensuremath{\Omega}}
\DeclareUnicodeCharacter{2200}{\ensuremath{\forall}}
\DeclareUnicodeCharacter{2203}{\ensuremath{\exists}}
\DeclareUnicodeCharacter{2207}{\ensuremath{\nabla}}
\DeclareUnicodeCharacter{2202}{\ensuremath{\partial}}
\DeclareUnicodeCharacter{2218}{\ensuremath{\circ}}
\DeclareUnicodeCharacter{2032}{\ensuremath{\prime}}
\DeclareUnicodeCharacter{2113}{\ensuremath{\ell}}
\DeclareUnicodeCharacter{210F}{\ensuremath{\hbar}}
\DeclareUnicodeCharacter{27E8}{\ensuremath{\langle}}
\DeclareUnicodeCharacter{27E9}{\ensuremath{\rangle}}
\DeclareUnicodeCharacter{2308}{\ensuremath{\lceil}}
\DeclareUnicodeCharacter{2309}{\ensuremath{\rceil}}
\DeclareUnicodeCharacter{230A}{\ensuremath{\lfloor}}
\DeclareUnicodeCharacter{230B}{\ensuremath{\rfloor}}
\DeclareUnicodeCharacter{222B}{\ensuremath{\int}}
\DeclareUnicodeCharacter{2211}{\ensuremath{\sum}}
\DeclareUnicodeCharacter{220F}{\ensuremath{\prod}}
\DeclareUnicodeCharacter{221E}{\ensuremath{\infty}}
\DeclareUnicodeCharacter{00B5}{\ensuremath{\mu}}
\DeclareUnicodeCharacter{2236}{\ensuremath{\colon}}
\DeclareUnicodeCharacter{2192}{\ensuremath{\to}}
\DeclareUnicodeCharacter{21A6}{\ensuremath{\mapsto}}
\DeclareUnicodeCharacter{21D2}{\ensuremath{\Rightarrow}}
\DeclareUnicodeCharacter{21D4}{\ensuremath{\Leftrightarrow}}
\DeclareUnicodeCharacter{2282}{\ensuremath{\subset}}
\DeclareUnicodeCharacter{2286}{\ensuremath{\subseteq}}
\DeclareUnicodeCharacter{2287}{\ensuremath{\supseteq}}
\DeclareUnicodeCharacter{222B}{\ensuremath{\int}}
\DeclareUnicodeCharacter{2205}{\ensuremath{\emptyset}}
\DeclareUnicodeCharacter{2299}{\ensuremath{\odot}}
\DeclareUnicodeCharacter{2220}{\ensuremath{\angle}}
\DeclareUnicodeCharacter{2135}{\ensuremath{\aleph}}
\DeclareUnicodeCharacter{207B}{\ensuremath{^{-}}}
\DeclareUnicodeCharacter{207A}{\ensuremath{^{+}}}
\DeclareUnicodeCharacter{220E}{\ensuremath{\blacksquare}}
\DeclareUnicodeCharacter{2061}{}
\DeclareUnicodeCharacter{2062}{}
\DeclareUnicodeCharacter{2063}{}
\DeclareUnicodeCharacter{2009}{\,}
\DeclareUnicodeCharacter{200A}{\,}
\DeclareUnicodeCharacter{200B}{}
\DeclareUnicodeCharacter{2212}{\ensuremath{-}}
\DeclareUnicodeCharacter{2194}{\ensuremath{\leftrightarrow}}
\DeclareUnicodeCharacter{2099}{\ensuremath{_{n}}}
\DeclareUnicodeCharacter{2098}{\ensuremath{_{m}}}
\DeclareUnicodeCharacter{2096}{\ensuremath{_{k}}}
\DeclareUnicodeCharacter{2095}{\ensuremath{_{h}}}
\DeclareUnicodeCharacter{2090}{\ensuremath{_{a}}}
\DeclareUnicodeCharacter{2091}{\ensuremath{_{e}}}
\DeclareUnicodeCharacter{2093}{\ensuremath{_{x}}}
\DeclareUnicodeCharacter{1D62}{\ensuremath{_{i}}}
\DeclareUnicodeCharacter{2C7C}{\ensuremath{_{j}}}
\DeclareUnicodeCharacter{2C7B}{\ensuremath{^{e}}}
\DeclareUnicodeCharacter{1D52}{\ensuremath{^{o}}}
\DeclareUnicodeCharacter{1D5}{\ensuremath{^{u}}}
\DeclareUnicodeCharacter{02E1}{\ensuremath{^{l}}}
\DeclareUnicodeCharacter{2295}{\ensuremath{\oplus}}
\DeclareUnicodeCharacter{27FA}{\ensuremath{\Longleftrightarrow}}
\DeclareUnicodeCharacter{27F9}{\ensuremath{\Longrightarrow}}
\DeclareUnicodeCharacter{2299}{\ensuremath{\odot}}
% Provide \tightlist used by pandoc-generated lists (no-op)
\providecommand{\tightlist}{\setlength{\itemsep}{0pt}\setlength{\parskip}{0pt}}
% pandoc longtable column specs: provide \real as passthrough and a
% `none` counter so that calc's \real{x} expansion path works.
\newcounter{none}
\providecommand{\real}[1]{#1}

% LISTINGS
\lstset{
    basicstyle=\small\ttfamily,
    breaklines=true, frame=single,
    keepspaces=true, showstringspaces=false,
    aboveskip=0.8em, belowskip=0.8em
}

% HEADER/FOOTER
\setlength{\headheight}{14pt}
\pagestyle{fancy}
\fancyhf{}
\fancyhead[R]{\thepage}
\renewcommand{\headrulewidth}{0.2pt}

% HYPERREF
\hypersetup{
    colorlinks=true,
    linkcolor=blue, citecolor=blue, urlcolor=blue,
    pdfauthor={Brieuc de La Fourniere}
}

% SPACING
\setstretch{1.2}
\setlength{\parskip}{0.4em}
\setlength{\parindent}{0pt}

% TITLE FORMATTING
\usepackage{titling}
\pretitle{\LARGE\bfseries}
\posttitle{\vspace{-0.4em}}
\preauthor{}
\postauthor{}
\predate{}
\postdate{}
\setlength{\droptitle}{-2.0em}

% CUSTOM COMMANDS — GIFT notation
\newcommand{\E}{\mathrm{E}}
\newcommand{\Gtwo}{\mathrm{G}_2}
\newcommand{\Kseven}{K_7}
\newcommand{\AdS}{\mathrm{AdS}}
\newcommand{\SU}{\mathrm{SU}}
\newcommand{\SO}{\mathrm{SO}}
\newcommand{\U}{\mathrm{U}}
\newcommand{\Sp}{\mathrm{Sp}}
\newcommand{\PSL}{\mathrm{PSL}}
\newcommand{\SM}{\mathrm{SM}}
\newcommand{\Tr}{\mathrm{Tr}}
\newcommand{\rank}{\mathrm{rank}}
\newcommand{\Vol}{\mathrm{Vol}}
\newcommand{\Hol}{\mathrm{Hol}}
\newcommand{\Ric}{\mathrm{Ric}}
\newcommand{\Rm}{\mathrm{Rm}}
\newcommand{\Scal}{\mathrm{Scal}}
\DeclareMathOperator{\diag}{diag}
\DeclareMathOperator{\cond}{cond}
\newcommand{\GIFT}{\textrm{GIFT}}
\newcommand{\CP}{\mathrm{CP}}
\newcommand{\CKM}{\mathrm{CKM}}
\newcommand{\PMNS}{\mathrm{PMNS}}
\newcommand{\EW}{\mathrm{EW}}
\newcommand{\Pl}{\mathrm{Pl}}
\newcommand{\DE}{\mathrm{DE}}
\newcommand{\DM}{\mathrm{DM}}
\newcommand{\KK}{\mathrm{KK}}
\newcommand{\NK}{\mathrm{NK}}
\newcommand{\TCS}{\mathrm{TCS}}
\newcommand{\PINN}{\mathrm{PINN}}
\newcommand{\btwo}{b_2}
\newcommand{\bthree}{b_3}
\newcommand{\typeI}{\textsc{I}}
\newcommand{\typeII}{\textsc{II}}
\newcommand{\typeIII}{\textsc{III}}
\newcommand{\typeIV}{\textsc{IV}}
\newcommand{\proven}{\textsc{Proven}}

% PDF bookmark safety
\pdfstringdefDisableCommands{%
  \def\textsubscript#1{#1}%
  \def\textsuperscript#1{#1}%
  \def\CP{CP}%
  \def\Gtwo{G2}%
  \def\Kseven{K7}%
  \def\GIFT{GIFT}%
  \def\E{E}%
  \def\SM{SM}%
}

% THEOREM ENVIRONMENTS
\newtheorem{theorem}{Theorem}[section]
\newtheorem{proposition}[theorem]{Proposition}
\newtheorem{lemma}[theorem]{Lemma}
\newtheorem{corollary}[theorem]{Corollary}
\theoremstyle{definition}
\newtheorem{definition}[theorem]{Definition}
\theoremstyle{remark}
\newtheorem{remark}[theorem]{Remark}

% pandoc >= 3.2.1 emits \pandocbounded around images; provide a fallback.
% Also clamp all images to \linewidth (natural-size figures overflowed the
% right margin -- the "truncated figures" flagged by the round-2 reviews).
% Kept here rather than in gift_preamble.tex, shared with the 3.4 build.
\providecommand{\pandocbounded}[1]{#1}
\makeatletter
\def\maxwidth{\ifdim\Gin@nat@width>\linewidth\linewidth\else\Gin@nat@width\fi}
\makeatother
\setkeys{Gin}{width=\maxwidth,keepaspectratio}
% md headings carry their own literal numbering ("3.2 What ..."); suppress
% the LaTeX section counters to avoid "0.2 1. Introduction" double numbering.
\setcounter{secnumdepth}{0}
% Loosen justification slightly: long technical identifiers (made breakable
% at build time via \allowbreak) still need extra stretch to avoid overfulls.
\emergencystretch=3em
\fancyhead[L]{The $K_7$ Framework v3.5: Supplement S4}
\hypersetup{pdftitle={The K7 Framework v3.5 Supplement S4: Sieve Diagnostics}}

\graphicspath{{../../../figures/}}

\title{%
\LARGE\textbf{Supplement S4: Sieve Diagnostics}\\[0.3em]
\large Coincidence-Probability Tables (Archived from v3.4 \S 7.5)%
}
\author{Brieuc de La Fourni\`ere \\ \textit{Independent researcher, Beaune, France}}
\date{July 2026 \quad\textbar\quad v3.5}

\begin{document}

\maketitle
\noindent\rule{\textwidth}{0.2pt}

\thispagestyle{fancy}

\emph{Companion to the K₇ framework paper v3.5. This supplement archives the coincidence-probability tables from the 3.4 §7.5 ``Coincidence Test'' section, which the 3.5 restructure demoted from statistical headline to internal-consistency diagnostics. The archived tables remain available for readers who want the previous framing. The 3.5 headline methodology is the Sieve reading (main paper §7); see the Arithmon methodology paper (Zenodo \href{https://doi.org/10.5281/zenodo.20666879}{10.5281/zenodo.20666879}) for the underlying construction.}

\section{S4.1 Purpose}\label{s4.1-purpose}

The 3.4 paper reported joint coincidence probabilities

\[
P_{\rm uniform} \sim 10^{-346}, \qquad P_{\rm algebraic} \sim 10^{-133}
\]

for the 33 Type I observables under (i) a uniform null model over \([0, 50\%]\) deviations, and (ii) an algebraic null model built from \(4.2 \cdot 10^6\) random formulas over the same 20 structural constants. These figures do not distinguish a \emph{genuine} survivor of a public sieve from a \emph{well-tuned} formula grammar; the 3.5 update supersedes them by the Sieve reading of §7 as headline, and preserves them here as diagnostic evidence of internal consistency.

\section{S4.2 Uniform Null Model (archived tables)}\label{s4.2-uniform-null-model-archived-tables}

Multiple statistical tests under the null hypothesis of uniform deviations in \([0, 50\%]\):

\begin{longtable}[]{@{}llll@{}}
\toprule\noalign{}
Test & Statistic & df & p-value \\
\midrule\noalign{}
\endhead
\bottomrule\noalign{}
\endlastfoot
\(\chi^2\) & 1063 & 33 & \(5.0 \times 10^{-202}\) \\
Fisher combined & 1100.5 & 66 & \(2.2 \times 10^{-187}\) \\
KS (pull normality) & 0.189 & n/a & 0.165 \\
\textbf{Combined uniform} & -- & -- & \textbf{\(10^{-346}\)} \\
\end{longtable}

\textbf{Pull distribution}: mean \(= -0.774\), std \(= 5.62\). 72.7\% of pulls within \(1\sigma\), 87.9\% within \(2\sigma\). The KS test (\(p = 0.165\)) is consistent with Gaussian pulls: the deviations are not cherry-picked.

\textbf{Reduced \(\chi^2 = 32.2\)} is large, driven by outliers (\(\delta_{\rm CP}\), \(\alpha_s(\text{RGE})\)). The Bayes factor \(\log_{10} = -2.3\) is inconclusive, reflecting this: the bulk of predictions match extremely well, but a few outliers inflate the tails. Removing the 4 known outliers gives reduced \(\chi^2 < 2\).

\section{S4.3 Algebraic Null Model (algebraic\_\allowbreak{}MC, archived)}\label{s4.3-algebraic-null-model-algebraic_mc-archived}

The uniform null model assumes predictions are random numbers. A stronger test: use random \emph{algebraic formulas} from the same 20 structural constants. \(4{,}188{,}086\) unique formula values were generated via exhaustive depth-1/2 enumeration (\(1.8 \cdot 10^6\)) plus \(3 \cdot 10^6\) random expression trees (depth 2--4), using 5 binary operations (\(+, -, \times, \div, \wedge\)) and 5 unary transforms (id, inv, \(\sqrt{}\), \(\square^2\), \(\ln\)).

\begin{longtable}[]{@{}llll@{}}
\toprule\noalign{}
Metric & K₇ framework & Random algebraic (median) & Factor \\
\midrule\noalign{}
\endhead
\bottomrule\noalign{}
\endlastfoot
Mean deviation & \textbf{0.73\%} & \(4.1 \cdot 10^9\) \% & \(5.6 \cdot 10^9\) \\
Exact matches (\textless{} 0.01\%) & \textbf{5} & 0 & -- \\
Within 1\% & \textbf{28} & 0 & -- \\
\end{longtable}

\textbf{Per-observable}: combined P(random algebraic formula matches the framework's precision) = \textbf{\(10^{-133}\)} (product over 33 observables). Only 0.02\% of \(4.2 \cdot 10^6\) formulas match \emph{any} observable within 0.01\%.

\textbf{Set-level}: 0 out of \(3{,}000{,}000\) random sets of 33 algebraic formulas achieved the framework's mean deviation, exact count, or within-1\% count. Under this null: \(P < 3.3 \cdot 10^{-7}\) for each metric.

\textbf{Interpretation (archived, 3.4 language).} The algebraic null model is far more generous than the uniform null (it generates formulas from the \emph{same constants}), yet the framework's performance remains unmatched across \(3 \cdot 10^6\) trials. Under both declared null models, chance agreement is excluded at extreme significance: \(P = 10^{-346}\) (uniform) and \(P = 10^{-133}\) (algebraic). These figures quantify agreement under each declared model; they do not constitute an absolute probability that the framework is correct.

\section{S4.4 What changed in 3.5 (Sieve reading)}\label{s4.4-what-changed-in-3.5-sieve-reading}

The Sieve reading of §7 in the main paper isolates individual \emph{survivors} rather than joint probabilities. Under the 4-null battery (uniform / algebraic / factorised / permuted) followed by budget-uniqueness ranking and Lean 4 R2 formal-identity flag, the outcome is:

\begin{itemize}
\tightlist
\item
  \(m_H / m_W\) (Route A): unique survivor at exact rank 1 budget-unique; carries the formal-identity flag via its Lean 4 R2 pre-registered identity.
\item
  Koide's \(Q = 2/3\): passes the battery with a narrow budget margin.
\item
  \(n_s\): passes the battery without a distinguishing budget signature.
\end{itemize}

The joint coincidence-probability language of §7.5 (this supplement) is not the right unit for the Sieve reading: two different jointness questions (probability under null vs certified survivor of a public sieve) are conflated in the 3.4 framing. The 3.5 restructure separates them and puts the certified-survivor language on the main line.

\section{S4.5 Reproducibility}\label{s4.5-reproducibility}

The archived null-model computations are reproducible from the same scripts used for the 3.4 release; the Sieve reading is reproducible via the Lean 4 module \texttt{Sieve/\allowbreak{}GStruct.\allowbreak{}lean} in the public \texttt{arithmon/\allowbreak{}Lean} repository. The Type I observable list is unchanged from 3.4 modulo the reconciliations documented in the CHANGELOG (observables reconciliation: 0.99\% Type I after the audit against the reconciled post-freeze sources used by the Arithmon/Sieve pipeline; the Type I tables of this paper use the earlier frozen PDG 2024 + NuFIT 6.0 dataset, mean 0.73\%, cf.~S2 §22; \(m_c/m_s\) erratum inlined).

\begin{center}\rule{\linewidth}{0.5pt}\end{center}

\emph{The K₇ Framework (formerly GIFT): Supplement S4: Sieve Diagnostics}


\bigskip
% ============================================
% GIFT v3.5 — Shared Bibliography
% ============================================
% Usage: % ============================================
% GIFT v3.5 — Shared Bibliography
% ============================================
% Usage: % ============================================
% GIFT v3.5 — Shared Bibliography
% ============================================
% Usage: \input{gift_3.5_bib} where the references section should appear.
% Sequential numbering [1]..[50]. Round-1 + round-2 IA-review updates applied
% (2026-07-07): [13] Kovalev 2005, [26] T2K/NOvA Nature 646,
% [27] NuFit 6.0 + 6.1, [45] Garbagnati-Sarti 2007, [46] LEPSUSYWG,
% [47] Fermilab DUNE 2031, [48] Chen-Hong (re-added, cited in S1), [49] Planck PR4, [50] CODATA.

\begin{thebibliography}{99}

\bibitem{1} Particle Data Group, ``Review of Particle Physics,'' \textit{Phys.\ Rev.\ D} \textbf{110}, 030001 (2024).

\bibitem{2} S.~Weinberg, ``Implications of dynamical symmetry breaking,'' \textit{Phys.\ Rev.\ D} \textbf{13}, 974 (1976).

\bibitem{3} Planck Collaboration, ``Planck 2018 results.\ VI.\ Cosmological parameters,'' \textit{Astron.\ Astrophys.} \textbf{641}, A6 (2020).

\bibitem{4} A.~G.~Riess et al., ``A comprehensive measurement of the local value of the Hubble constant,'' \textit{ApJL} \textbf{934}, L7 (2022).

\bibitem{5} C.~D.~Froggatt, H.~B.~Nielsen, ``Hierarchy of quark masses, Cabibbo angles and CP violation,'' \textit{Nucl.\ Phys.\ B} \textbf{147}, 277 (1979).

\bibitem{6} Y.~Koide, ``Fermion-boson two-body model of quarks and leptons and Cabibbo mixing,'' \textit{Lett.\ Nuovo Cim.} \textbf{34}, 201 (1982).

\bibitem{7} C.~Furey, \textit{Standard Model Physics from an Algebra?}, PhD thesis, University of Waterloo (2015).

\bibitem{8} N.~Furey, M.~J.~Hughes, ``One generation of standard model Weyl representations as a single copy of $\mathbb{R} \otimes \mathbb{C} \otimes \mathbb{H} \otimes \mathbb{O}$,'' \textit{Phys.\ Lett.\ B} \textbf{831}, 137186 (2022).

\bibitem{9} R.~Wilson, ``$E_8$ and Standard Model plus gravity,'' arXiv:2404.18938 [hep-th] (2024).

\bibitem{10} T.~P.~Singh et al., ``An $E_8 \otimes E_8$ unification of the Standard Model with pre-gravitation,'' arXiv:2206.06911v3 (2024).

\bibitem{11} B.~S.~Acharya, S.~Gukov, ``M-theory and singularities of exceptional holonomy manifolds,'' \textit{Phys.\ Rep.} \textbf{392}, 121 (2004).

\bibitem{12} L.~Foscolo et al., ``Complete non-compact $G_2$-manifolds from asymptotically conical Calabi--Yau 3-folds,'' \textit{Duke Math.\ J.} \textbf{170}, 3 (2021).

\bibitem{13} A.~Kovalev, ``Coassociative K3 fibrations of compact $G_2$-manifolds,'' arXiv:math/0511150 (2005).

\bibitem{14} M.~Haskins, J.~Nordstr\"om, ``Extra-twisted connected sum $G_2$-manifolds,'' arXiv:1809.09083 (2022); \textit{Ann.\ Glob.\ Anal.\ Geom.} \textbf{64}, art.\ 2 (2023).

\bibitem{15} G.~Bera, ``Associative submanifolds in twisted connected sum $G_2$-manifolds,'' arXiv:2209.00156 (2022).

\bibitem{16} D.~Crowley, S.~Goette, J.~Nordstr\"om, ``An analytic invariant of $G_2$ manifolds,'' \textit{Invent.\ Math.} \textbf{239}(3), 865--907 (2025).

\bibitem{17} T.~Dray, C.~A.~Manogue, \textit{The Geometry of the Octonions}, World Scientific (2015).

\bibitem{18} J.~F.~Adams, \textit{Lectures on Exceptional Lie Groups}, University of Chicago Press (1996).

\bibitem{19} D.~J.~Gross et al., ``Heterotic string theory,'' \textit{Nucl.\ Phys.\ B} \textbf{256}, 253 (1985).

\bibitem{20} D.~D.~Joyce, \textit{Compact Manifolds with Special Holonomy}, Oxford University Press (2000).

\bibitem{21} A.~Kovalev, ``Twisted connected sums and special Riemannian holonomy,'' \textit{J.\ Reine Angew.\ Math.} \textbf{565}, 125 (2003).

\bibitem{22} A.~Corti et al., ``$G_2$-manifolds and associative submanifolds via semi-Fano 3-folds,'' \textit{Duke Math.\ J.} \textbf{164}, 1971 (2015).

\bibitem{23} R.~Harvey, H.~B.~Lawson, ``Calibrated geometries,'' \textit{Acta Math.} \textbf{148}, 47 (1982).

\bibitem{24} B.~S.~Acharya, ``M-theory, $G_2$-manifolds and four-dimensional physics,'' \textit{Class.\ Quant.\ Grav.} \textbf{19}, 5619 (2002).

\bibitem{25} B.~S.~Acharya, E.~Witten, ``Chiral fermions from manifolds of $G_2$ holonomy,'' arXiv:hep-th/0109152 (2001).

\bibitem{26} T2K and NOvA Collaborations, ``Joint analysis of long-baseline neutrino oscillations,'' \textit{Nature} \textbf{646}, 818--824 (2025), \href{https://doi.org/10.1038/s41586-025-09599-3}{doi:10.1038/s41586-025-09599-3}.

\bibitem{27} I.~Esteban et al.\ (NuFIT), ``NuFit-6.0: three-flavour global analysis of neutrino oscillation experiments,'' \textit{JHEP} \textbf{12} (2024) 216, arXiv:2410.05380; NuFIT 6.1 web release, \url{https://www.nu-fit.org} (2025).

\bibitem{28} L.~de Moura, S.~Ullrich, ``The Lean 4 theorem prover and programming language,'' in \textit{CADE 28}, p.\ 625 (2021).

\bibitem{29} mathlib Community, \textit{mathlib4}, \url{https://github.com/leanprover-community/mathlib4}.

\bibitem{30} DUNE Collaboration, ``Long-baseline neutrino facility (LBNF) and DUNE conceptual design report,'' FERMILAB-TM-2696 (2020).

\bibitem{31} DUNE Collaboration, ``Deep Underground Neutrino Experiment (DUNE) far detector technical design report,'' arXiv:2103.04797 (2021).

\bibitem{32} E.~Witten, ``Strong coupling expansion of Calabi-Yau compactification,'' \textit{Nucl.\ Phys.\ B} \textbf{471}, 135 (1996).

\bibitem{33} B.~S.~Acharya et al., ``M-theory solution to the hierarchy problem,'' \textit{Phys.\ Rev.\ D} \textbf{76}, 126010 (2007).

\bibitem{34} M.~Atiyah, E.~Witten, ``M-theory dynamics on a manifold of $G_2$-holonomy,'' \textit{Adv.\ Theor.\ Math.\ Phys.} \textbf{6}, 1 (2002).

\bibitem{35} G.~Kane, \textit{String Theory and the Real World} (2017).

\bibitem{36} J.~Distler, S.~Garibaldi, ``There is no `Theory of Everything' inside $E_8$,'' \textit{Commun.\ Math.\ Phys.} \textbf{298}, 419 (2010).

\bibitem{37} J.~C.~Baez, ``Octonions and the Standard Model,'' \url{https://math.ucr.edu/home/baez/standard/} (2020--2025).

\bibitem{38} Springer Nature, ``Artificial intelligence (AI) policy,'' \url{https://www.springernature.com/gp/policies} (2024).

\bibitem{39} J.~A.~Wheeler, ``Information, physics, quantum: the search for links,'' in \textit{Complexity, Entropy, and the Physics of Information}, W.~H.~Zurek (ed.), Addison-Wesley (1990), pp.~3--28.

\bibitem{40} J.~Worrall, ``Structural realism: the best of both worlds?,'' \textit{Dialectica} \textbf{43}, 99--124 (1989).

\bibitem{41} J.~Ladyman, D.~Ross, \textit{Every Thing Must Go: Metaphysics Naturalized}, Oxford University Press (2007).

\bibitem{42} R.~Pin\v{c}\'ak, A.~Pigazzini, M.~Pudl\'ak, E.~Barto\v{s}, ``Geometric origin of a stable black hole remnant from torsion in $G_2$-manifold geometry,'' \textit{Gen.\ Rel.\ Grav.} \textbf{58}, 29 (2026), \href{https://doi.org/10.1007/s10714-026-03528-z}{doi:10.1007/s10714-026-03528-z}.

\bibitem{43} D.~Joyce, S.~Karigiannis, ``A new construction of compact torsion-free $G_2$-manifolds by gluing families of Eguchi--Hanson spaces,'' \textit{J.\ Differential Geom.} (2021), arXiv:1707.09325 (2017).

\bibitem{44} S.~Mukai, ``Finite groups of automorphisms of K3 surfaces and the Mathieu group,'' \textit{Invent.\ Math.} \textbf{94}, 183--221 (1988).

\bibitem{45} A.~Garbagnati, A.~Sarti, ``Symplectic automorphisms of prime order on K3 surfaces,'' \textit{J.\ Algebra} \textbf{318}, 323--350 (2007), arXiv:math/0603742.

\bibitem{46} LEP SUSY Working Group (ALEPH, DELPHI, L3, OPAL), ``Combined LEP chargino results, up to 208 GeV for large $m_0$,'' LEPSUSYWG/01-03.1 (2001), \url{https://lepsusy.web.cern.ch/lepsusy/www/inos_moriond01/charginos_pub.html}.

\bibitem{47} Fermilab, ``DUNE first-beam target of 2031: LBNF/DUNE-US milestone communication'' (2026-05-07), \url{https://news.fnal.gov/}.

\bibitem{48} J.~Chen, H.~Hong, ``Intermediate curvature and splitting theorem,'' arXiv:2604.26529 (2026).

\bibitem{49} M.~Tristram et al., ``Cosmological parameters derived from the final Planck data release (PR4),'' \textit{Astron.\ Astrophys.} \textbf{682}, A37 (2024).

\bibitem{50} E.~Tiesinga, P.~J.~Mohr, D.~B.~Newell, B.~N.~Taylor, ``CODATA recommended values of the fundamental physical constants: 2022,'' \textit{Rev.\ Mod.\ Phys.} \textbf{97}, 025002 (2025), \url{https://physics.nist.gov/cuu/Constants}.

\bibitem{51} V.~V.~Nikulin, ``Integer symmetric bilinear forms and some of their geometric applications,'' \textit{Izv.\ Akad.\ Nauk SSSR Ser.\ Mat.} \textbf{43} (1979); English transl.\ \textit{Math.\ USSR Izv.} \textbf{14} (1980).

\end{thebibliography}
 where the references section should appear.
% Sequential numbering [1]..[50]. Round-1 + round-2 IA-review updates applied
% (2026-07-07): [13] Kovalev 2005, [26] T2K/NOvA Nature 646,
% [27] NuFit 6.0 + 6.1, [45] Garbagnati-Sarti 2007, [46] LEPSUSYWG,
% [47] Fermilab DUNE 2031, [48] Chen-Hong (re-added, cited in S1), [49] Planck PR4, [50] CODATA.

\begin{thebibliography}{99}

\bibitem{1} Particle Data Group, ``Review of Particle Physics,'' \textit{Phys.\ Rev.\ D} \textbf{110}, 030001 (2024).

\bibitem{2} S.~Weinberg, ``Implications of dynamical symmetry breaking,'' \textit{Phys.\ Rev.\ D} \textbf{13}, 974 (1976).

\bibitem{3} Planck Collaboration, ``Planck 2018 results.\ VI.\ Cosmological parameters,'' \textit{Astron.\ Astrophys.} \textbf{641}, A6 (2020).

\bibitem{4} A.~G.~Riess et al., ``A comprehensive measurement of the local value of the Hubble constant,'' \textit{ApJL} \textbf{934}, L7 (2022).

\bibitem{5} C.~D.~Froggatt, H.~B.~Nielsen, ``Hierarchy of quark masses, Cabibbo angles and CP violation,'' \textit{Nucl.\ Phys.\ B} \textbf{147}, 277 (1979).

\bibitem{6} Y.~Koide, ``Fermion-boson two-body model of quarks and leptons and Cabibbo mixing,'' \textit{Lett.\ Nuovo Cim.} \textbf{34}, 201 (1982).

\bibitem{7} C.~Furey, \textit{Standard Model Physics from an Algebra?}, PhD thesis, University of Waterloo (2015).

\bibitem{8} N.~Furey, M.~J.~Hughes, ``One generation of standard model Weyl representations as a single copy of $\mathbb{R} \otimes \mathbb{C} \otimes \mathbb{H} \otimes \mathbb{O}$,'' \textit{Phys.\ Lett.\ B} \textbf{831}, 137186 (2022).

\bibitem{9} R.~Wilson, ``$E_8$ and Standard Model plus gravity,'' arXiv:2404.18938 [hep-th] (2024).

\bibitem{10} T.~P.~Singh et al., ``An $E_8 \otimes E_8$ unification of the Standard Model with pre-gravitation,'' arXiv:2206.06911v3 (2024).

\bibitem{11} B.~S.~Acharya, S.~Gukov, ``M-theory and singularities of exceptional holonomy manifolds,'' \textit{Phys.\ Rep.} \textbf{392}, 121 (2004).

\bibitem{12} L.~Foscolo et al., ``Complete non-compact $G_2$-manifolds from asymptotically conical Calabi--Yau 3-folds,'' \textit{Duke Math.\ J.} \textbf{170}, 3 (2021).

\bibitem{13} A.~Kovalev, ``Coassociative K3 fibrations of compact $G_2$-manifolds,'' arXiv:math/0511150 (2005).

\bibitem{14} M.~Haskins, J.~Nordstr\"om, ``Extra-twisted connected sum $G_2$-manifolds,'' arXiv:1809.09083 (2022); \textit{Ann.\ Glob.\ Anal.\ Geom.} \textbf{64}, art.\ 2 (2023).

\bibitem{15} G.~Bera, ``Associative submanifolds in twisted connected sum $G_2$-manifolds,'' arXiv:2209.00156 (2022).

\bibitem{16} D.~Crowley, S.~Goette, J.~Nordstr\"om, ``An analytic invariant of $G_2$ manifolds,'' \textit{Invent.\ Math.} \textbf{239}(3), 865--907 (2025).

\bibitem{17} T.~Dray, C.~A.~Manogue, \textit{The Geometry of the Octonions}, World Scientific (2015).

\bibitem{18} J.~F.~Adams, \textit{Lectures on Exceptional Lie Groups}, University of Chicago Press (1996).

\bibitem{19} D.~J.~Gross et al., ``Heterotic string theory,'' \textit{Nucl.\ Phys.\ B} \textbf{256}, 253 (1985).

\bibitem{20} D.~D.~Joyce, \textit{Compact Manifolds with Special Holonomy}, Oxford University Press (2000).

\bibitem{21} A.~Kovalev, ``Twisted connected sums and special Riemannian holonomy,'' \textit{J.\ Reine Angew.\ Math.} \textbf{565}, 125 (2003).

\bibitem{22} A.~Corti et al., ``$G_2$-manifolds and associative submanifolds via semi-Fano 3-folds,'' \textit{Duke Math.\ J.} \textbf{164}, 1971 (2015).

\bibitem{23} R.~Harvey, H.~B.~Lawson, ``Calibrated geometries,'' \textit{Acta Math.} \textbf{148}, 47 (1982).

\bibitem{24} B.~S.~Acharya, ``M-theory, $G_2$-manifolds and four-dimensional physics,'' \textit{Class.\ Quant.\ Grav.} \textbf{19}, 5619 (2002).

\bibitem{25} B.~S.~Acharya, E.~Witten, ``Chiral fermions from manifolds of $G_2$ holonomy,'' arXiv:hep-th/0109152 (2001).

\bibitem{26} T2K and NOvA Collaborations, ``Joint analysis of long-baseline neutrino oscillations,'' \textit{Nature} \textbf{646}, 818--824 (2025), \href{https://doi.org/10.1038/s41586-025-09599-3}{doi:10.1038/s41586-025-09599-3}.

\bibitem{27} I.~Esteban et al.\ (NuFIT), ``NuFit-6.0: three-flavour global analysis of neutrino oscillation experiments,'' \textit{JHEP} \textbf{12} (2024) 216, arXiv:2410.05380; NuFIT 6.1 web release, \url{https://www.nu-fit.org} (2025).

\bibitem{28} L.~de Moura, S.~Ullrich, ``The Lean 4 theorem prover and programming language,'' in \textit{CADE 28}, p.\ 625 (2021).

\bibitem{29} mathlib Community, \textit{mathlib4}, \url{https://github.com/leanprover-community/mathlib4}.

\bibitem{30} DUNE Collaboration, ``Long-baseline neutrino facility (LBNF) and DUNE conceptual design report,'' FERMILAB-TM-2696 (2020).

\bibitem{31} DUNE Collaboration, ``Deep Underground Neutrino Experiment (DUNE) far detector technical design report,'' arXiv:2103.04797 (2021).

\bibitem{32} E.~Witten, ``Strong coupling expansion of Calabi-Yau compactification,'' \textit{Nucl.\ Phys.\ B} \textbf{471}, 135 (1996).

\bibitem{33} B.~S.~Acharya et al., ``M-theory solution to the hierarchy problem,'' \textit{Phys.\ Rev.\ D} \textbf{76}, 126010 (2007).

\bibitem{34} M.~Atiyah, E.~Witten, ``M-theory dynamics on a manifold of $G_2$-holonomy,'' \textit{Adv.\ Theor.\ Math.\ Phys.} \textbf{6}, 1 (2002).

\bibitem{35} G.~Kane, \textit{String Theory and the Real World} (2017).

\bibitem{36} J.~Distler, S.~Garibaldi, ``There is no `Theory of Everything' inside $E_8$,'' \textit{Commun.\ Math.\ Phys.} \textbf{298}, 419 (2010).

\bibitem{37} J.~C.~Baez, ``Octonions and the Standard Model,'' \url{https://math.ucr.edu/home/baez/standard/} (2020--2025).

\bibitem{38} Springer Nature, ``Artificial intelligence (AI) policy,'' \url{https://www.springernature.com/gp/policies} (2024).

\bibitem{39} J.~A.~Wheeler, ``Information, physics, quantum: the search for links,'' in \textit{Complexity, Entropy, and the Physics of Information}, W.~H.~Zurek (ed.), Addison-Wesley (1990), pp.~3--28.

\bibitem{40} J.~Worrall, ``Structural realism: the best of both worlds?,'' \textit{Dialectica} \textbf{43}, 99--124 (1989).

\bibitem{41} J.~Ladyman, D.~Ross, \textit{Every Thing Must Go: Metaphysics Naturalized}, Oxford University Press (2007).

\bibitem{42} R.~Pin\v{c}\'ak, A.~Pigazzini, M.~Pudl\'ak, E.~Barto\v{s}, ``Geometric origin of a stable black hole remnant from torsion in $G_2$-manifold geometry,'' \textit{Gen.\ Rel.\ Grav.} \textbf{58}, 29 (2026), \href{https://doi.org/10.1007/s10714-026-03528-z}{doi:10.1007/s10714-026-03528-z}.

\bibitem{43} D.~Joyce, S.~Karigiannis, ``A new construction of compact torsion-free $G_2$-manifolds by gluing families of Eguchi--Hanson spaces,'' \textit{J.\ Differential Geom.} (2021), arXiv:1707.09325 (2017).

\bibitem{44} S.~Mukai, ``Finite groups of automorphisms of K3 surfaces and the Mathieu group,'' \textit{Invent.\ Math.} \textbf{94}, 183--221 (1988).

\bibitem{45} A.~Garbagnati, A.~Sarti, ``Symplectic automorphisms of prime order on K3 surfaces,'' \textit{J.\ Algebra} \textbf{318}, 323--350 (2007), arXiv:math/0603742.

\bibitem{46} LEP SUSY Working Group (ALEPH, DELPHI, L3, OPAL), ``Combined LEP chargino results, up to 208 GeV for large $m_0$,'' LEPSUSYWG/01-03.1 (2001), \url{https://lepsusy.web.cern.ch/lepsusy/www/inos_moriond01/charginos_pub.html}.

\bibitem{47} Fermilab, ``DUNE first-beam target of 2031: LBNF/DUNE-US milestone communication'' (2026-05-07), \url{https://news.fnal.gov/}.

\bibitem{48} J.~Chen, H.~Hong, ``Intermediate curvature and splitting theorem,'' arXiv:2604.26529 (2026).

\bibitem{49} M.~Tristram et al., ``Cosmological parameters derived from the final Planck data release (PR4),'' \textit{Astron.\ Astrophys.} \textbf{682}, A37 (2024).

\bibitem{50} E.~Tiesinga, P.~J.~Mohr, D.~B.~Newell, B.~N.~Taylor, ``CODATA recommended values of the fundamental physical constants: 2022,'' \textit{Rev.\ Mod.\ Phys.} \textbf{97}, 025002 (2025), \url{https://physics.nist.gov/cuu/Constants}.

\bibitem{51} V.~V.~Nikulin, ``Integer symmetric bilinear forms and some of their geometric applications,'' \textit{Izv.\ Akad.\ Nauk SSSR Ser.\ Mat.} \textbf{43} (1979); English transl.\ \textit{Math.\ USSR Izv.} \textbf{14} (1980).

\end{thebibliography}
 where the references section should appear.
% Sequential numbering [1]..[50]. Round-1 + round-2 IA-review updates applied
% (2026-07-07): [13] Kovalev 2005, [26] T2K/NOvA Nature 646,
% [27] NuFit 6.0 + 6.1, [45] Garbagnati-Sarti 2007, [46] LEPSUSYWG,
% [47] Fermilab DUNE 2031, [48] Chen-Hong (re-added, cited in S1), [49] Planck PR4, [50] CODATA.

\begin{thebibliography}{99}

\bibitem{1} Particle Data Group, ``Review of Particle Physics,'' \textit{Phys.\ Rev.\ D} \textbf{110}, 030001 (2024).

\bibitem{2} S.~Weinberg, ``Implications of dynamical symmetry breaking,'' \textit{Phys.\ Rev.\ D} \textbf{13}, 974 (1976).

\bibitem{3} Planck Collaboration, ``Planck 2018 results.\ VI.\ Cosmological parameters,'' \textit{Astron.\ Astrophys.} \textbf{641}, A6 (2020).

\bibitem{4} A.~G.~Riess et al., ``A comprehensive measurement of the local value of the Hubble constant,'' \textit{ApJL} \textbf{934}, L7 (2022).

\bibitem{5} C.~D.~Froggatt, H.~B.~Nielsen, ``Hierarchy of quark masses, Cabibbo angles and CP violation,'' \textit{Nucl.\ Phys.\ B} \textbf{147}, 277 (1979).

\bibitem{6} Y.~Koide, ``Fermion-boson two-body model of quarks and leptons and Cabibbo mixing,'' \textit{Lett.\ Nuovo Cim.} \textbf{34}, 201 (1982).

\bibitem{7} C.~Furey, \textit{Standard Model Physics from an Algebra?}, PhD thesis, University of Waterloo (2015).

\bibitem{8} N.~Furey, M.~J.~Hughes, ``One generation of standard model Weyl representations as a single copy of $\mathbb{R} \otimes \mathbb{C} \otimes \mathbb{H} \otimes \mathbb{O}$,'' \textit{Phys.\ Lett.\ B} \textbf{831}, 137186 (2022).

\bibitem{9} R.~Wilson, ``$E_8$ and Standard Model plus gravity,'' arXiv:2404.18938 [hep-th] (2024).

\bibitem{10} T.~P.~Singh et al., ``An $E_8 \otimes E_8$ unification of the Standard Model with pre-gravitation,'' arXiv:2206.06911v3 (2024).

\bibitem{11} B.~S.~Acharya, S.~Gukov, ``M-theory and singularities of exceptional holonomy manifolds,'' \textit{Phys.\ Rep.} \textbf{392}, 121 (2004).

\bibitem{12} L.~Foscolo et al., ``Complete non-compact $G_2$-manifolds from asymptotically conical Calabi--Yau 3-folds,'' \textit{Duke Math.\ J.} \textbf{170}, 3 (2021).

\bibitem{13} A.~Kovalev, ``Coassociative K3 fibrations of compact $G_2$-manifolds,'' arXiv:math/0511150 (2005).

\bibitem{14} M.~Haskins, J.~Nordstr\"om, ``Extra-twisted connected sum $G_2$-manifolds,'' arXiv:1809.09083 (2022); \textit{Ann.\ Glob.\ Anal.\ Geom.} \textbf{64}, art.\ 2 (2023).

\bibitem{15} G.~Bera, ``Associative submanifolds in twisted connected sum $G_2$-manifolds,'' arXiv:2209.00156 (2022).

\bibitem{16} D.~Crowley, S.~Goette, J.~Nordstr\"om, ``An analytic invariant of $G_2$ manifolds,'' \textit{Invent.\ Math.} \textbf{239}(3), 865--907 (2025).

\bibitem{17} T.~Dray, C.~A.~Manogue, \textit{The Geometry of the Octonions}, World Scientific (2015).

\bibitem{18} J.~F.~Adams, \textit{Lectures on Exceptional Lie Groups}, University of Chicago Press (1996).

\bibitem{19} D.~J.~Gross et al., ``Heterotic string theory,'' \textit{Nucl.\ Phys.\ B} \textbf{256}, 253 (1985).

\bibitem{20} D.~D.~Joyce, \textit{Compact Manifolds with Special Holonomy}, Oxford University Press (2000).

\bibitem{21} A.~Kovalev, ``Twisted connected sums and special Riemannian holonomy,'' \textit{J.\ Reine Angew.\ Math.} \textbf{565}, 125 (2003).

\bibitem{22} A.~Corti et al., ``$G_2$-manifolds and associative submanifolds via semi-Fano 3-folds,'' \textit{Duke Math.\ J.} \textbf{164}, 1971 (2015).

\bibitem{23} R.~Harvey, H.~B.~Lawson, ``Calibrated geometries,'' \textit{Acta Math.} \textbf{148}, 47 (1982).

\bibitem{24} B.~S.~Acharya, ``M-theory, $G_2$-manifolds and four-dimensional physics,'' \textit{Class.\ Quant.\ Grav.} \textbf{19}, 5619 (2002).

\bibitem{25} B.~S.~Acharya, E.~Witten, ``Chiral fermions from manifolds of $G_2$ holonomy,'' arXiv:hep-th/0109152 (2001).

\bibitem{26} T2K and NOvA Collaborations, ``Joint analysis of long-baseline neutrino oscillations,'' \textit{Nature} \textbf{646}, 818--824 (2025), \href{https://doi.org/10.1038/s41586-025-09599-3}{doi:10.1038/s41586-025-09599-3}.

\bibitem{27} I.~Esteban et al.\ (NuFIT), ``NuFit-6.0: three-flavour global analysis of neutrino oscillation experiments,'' \textit{JHEP} \textbf{12} (2024) 216, arXiv:2410.05380; NuFIT 6.1 web release, \url{https://www.nu-fit.org} (2025).

\bibitem{28} L.~de Moura, S.~Ullrich, ``The Lean 4 theorem prover and programming language,'' in \textit{CADE 28}, p.\ 625 (2021).

\bibitem{29} mathlib Community, \textit{mathlib4}, \url{https://github.com/leanprover-community/mathlib4}.

\bibitem{30} DUNE Collaboration, ``Long-baseline neutrino facility (LBNF) and DUNE conceptual design report,'' FERMILAB-TM-2696 (2020).

\bibitem{31} DUNE Collaboration, ``Deep Underground Neutrino Experiment (DUNE) far detector technical design report,'' arXiv:2103.04797 (2021).

\bibitem{32} E.~Witten, ``Strong coupling expansion of Calabi-Yau compactification,'' \textit{Nucl.\ Phys.\ B} \textbf{471}, 135 (1996).

\bibitem{33} B.~S.~Acharya et al., ``M-theory solution to the hierarchy problem,'' \textit{Phys.\ Rev.\ D} \textbf{76}, 126010 (2007).

\bibitem{34} M.~Atiyah, E.~Witten, ``M-theory dynamics on a manifold of $G_2$-holonomy,'' \textit{Adv.\ Theor.\ Math.\ Phys.} \textbf{6}, 1 (2002).

\bibitem{35} G.~Kane, \textit{String Theory and the Real World} (2017).

\bibitem{36} J.~Distler, S.~Garibaldi, ``There is no `Theory of Everything' inside $E_8$,'' \textit{Commun.\ Math.\ Phys.} \textbf{298}, 419 (2010).

\bibitem{37} J.~C.~Baez, ``Octonions and the Standard Model,'' \url{https://math.ucr.edu/home/baez/standard/} (2020--2025).

\bibitem{38} Springer Nature, ``Artificial intelligence (AI) policy,'' \url{https://www.springernature.com/gp/policies} (2024).

\bibitem{39} J.~A.~Wheeler, ``Information, physics, quantum: the search for links,'' in \textit{Complexity, Entropy, and the Physics of Information}, W.~H.~Zurek (ed.), Addison-Wesley (1990), pp.~3--28.

\bibitem{40} J.~Worrall, ``Structural realism: the best of both worlds?,'' \textit{Dialectica} \textbf{43}, 99--124 (1989).

\bibitem{41} J.~Ladyman, D.~Ross, \textit{Every Thing Must Go: Metaphysics Naturalized}, Oxford University Press (2007).

\bibitem{42} R.~Pin\v{c}\'ak, A.~Pigazzini, M.~Pudl\'ak, E.~Barto\v{s}, ``Geometric origin of a stable black hole remnant from torsion in $G_2$-manifold geometry,'' \textit{Gen.\ Rel.\ Grav.} \textbf{58}, 29 (2026), \href{https://doi.org/10.1007/s10714-026-03528-z}{doi:10.1007/s10714-026-03528-z}.

\bibitem{43} D.~Joyce, S.~Karigiannis, ``A new construction of compact torsion-free $G_2$-manifolds by gluing families of Eguchi--Hanson spaces,'' \textit{J.\ Differential Geom.} (2021), arXiv:1707.09325 (2017).

\bibitem{44} S.~Mukai, ``Finite groups of automorphisms of K3 surfaces and the Mathieu group,'' \textit{Invent.\ Math.} \textbf{94}, 183--221 (1988).

\bibitem{45} A.~Garbagnati, A.~Sarti, ``Symplectic automorphisms of prime order on K3 surfaces,'' \textit{J.\ Algebra} \textbf{318}, 323--350 (2007), arXiv:math/0603742.

\bibitem{46} LEP SUSY Working Group (ALEPH, DELPHI, L3, OPAL), ``Combined LEP chargino results, up to 208 GeV for large $m_0$,'' LEPSUSYWG/01-03.1 (2001), \url{https://lepsusy.web.cern.ch/lepsusy/www/inos_moriond01/charginos_pub.html}.

\bibitem{47} Fermilab, ``DUNE first-beam target of 2031: LBNF/DUNE-US milestone communication'' (2026-05-07), \url{https://news.fnal.gov/}.

\bibitem{48} J.~Chen, H.~Hong, ``Intermediate curvature and splitting theorem,'' arXiv:2604.26529 (2026).

\bibitem{49} M.~Tristram et al., ``Cosmological parameters derived from the final Planck data release (PR4),'' \textit{Astron.\ Astrophys.} \textbf{682}, A37 (2024).

\bibitem{50} E.~Tiesinga, P.~J.~Mohr, D.~B.~Newell, B.~N.~Taylor, ``CODATA recommended values of the fundamental physical constants: 2022,'' \textit{Rev.\ Mod.\ Phys.} \textbf{97}, 025002 (2025), \url{https://physics.nist.gov/cuu/Constants}.

\bibitem{51} V.~V.~Nikulin, ``Integer symmetric bilinear forms and some of their geometric applications,'' \textit{Izv.\ Akad.\ Nauk SSSR Ser.\ Mat.} \textbf{43} (1979); English transl.\ \textit{Math.\ USSR Izv.} \textbf{14} (1980).

\end{thebibliography}


\end{document}
